\lyseis{6.3}{Ανάλογα ποσά}
\sol{1} $y=3x$: για $x=0,\,1,\,2,\,5,\,0{,}4$ προκύπτει $y=0,\,3,\,6,\,15,\,1{,}2$. Για $y=9$: $x=9:3=3$. Για $y=4{,}5$: $x=1{,}5$.
\sol{2} $y=\frac{2}{3}x$: για $x=0,\,3,\,6,\,1$ είναι $y=0,\,2,\,4,\,\frac{2}{3}$. Για $y=10$: $x=10:\frac{2}{3}=15$.
\sol{3} α) Ναι: $\frac{5}{2}=\frac{10}{4}=\frac{15}{6}=\frac{20}{8}=2{,}5$. β) Όχι: $\frac{2}{1}=2$, $\frac{3}{2}=1{,}5$. γ) Ναι: $\alpha=\frac{1}{3}$. δ) Όχι: $\frac{5}{1}=\frac{10}{2}=5$ αλλά $\frac{21}{4}=5{,}25$.
\sol{4} α) $y=1{,}5x$, ανάλογα. β) $y=4x$, ανάλογα. γ) $y=x\cdot x$, όχι ανάλογα (π.χ. $x=1\to1$, $x=2\to4$, όχι $2$).
\sol{5} $\alpha=\frac{15}{5}=3$, και $\frac{24}{8}=3$, $\frac{36}{12}=3$.
\sol{6} α) Σωστό. β) Σωστό. γ) Λάθος: για $x=0$ δίνει $y=2$. δ) Σωστό. ε) Λάθος: $y=x^2$.
\sol{7} $0{,}5\to4$, $1\to8$, $1{,}5\to12$, $2{,}5\to20$, $4\to32$ €. Με $52$ €: $52:8=6{,}5$ κιλά.
\sol{8} $20\cdot1{,}85=37$ € και $35\cdot1{,}85=64{,}75$ €. Με $50$ €: $50:1{,}85\approx27{,}03$ λίτρα.
\sol{9} $2{,}5\cdot4{,}8=12$ €, $0{,}75\cdot4{,}8=3{,}6$ €. Με $33{,}6$ €: $33{,}6:4{,}8=7$ μέτρα.
\sol{10} $30$ g αλάτι και $75$ g διαλύματος.
\sol{11} $0{,}24\cdot50=12$ €, $0{,}24\cdot200=48$ €. Για ΦΠΑ $30$ €: αξία $30:0{,}24=125$ €.
\sol{12} $9$ km, $36$ km, $63$ km. Για $27$ km: $27:18=1{,}5$ ώρες.
\sol{13} $2{,}5\cdot8=20$ m$^2$. Για $60$ m$^2$: $60:8=7{,}5$ λίτρα.
\sol{14} Το ύψος και η ηλικία δεν είναι ανάλογα: ο λόγος $\frac{\text{ύψος}}{\text{ηλικία}}$ δεν μένει σταθερός ($\frac{1{,}28}{8}=0{,}16$ αλλά $\frac{1{,}55}{12}\approx0{,}13$). Ο άνθρωπος σταματά να μεγαλώνει περίπου στα $18$--$20$.
\sol{15} Ανάλογα: α, γ, ε. Όχι: β, δ.
\sol{16} $\frac{y}{z}=\frac{2x}{5x}=\frac{2}{5}$, σταθερός, άρα ανάλογα με $\alpha=\frac{2}{5}$ (και $z=\frac{5}{2}y$).
\sol{17} $y=3\cdot7=21$. Δηλαδή $2+5=7$ και $6+15=21$.
\sol{18} Ελεύθερη απάντηση.
\sol{19} Ελεύθερη απάντηση. Π.χ. ανάλογα: τα κιλά και η τιμή. Όχι ανάλογα: ηλικία και ύψος.
